% ===========================================================================
% Verified Twin Execution Kernel -- correctness report
%
% Build:  make -C proof        (latexmk; output proof/build/main.pdf)
%
% This is a rigorous pen-and-paper argument. It has NOT been mechanically
% verified; mechanisation (Lean/Coq/Isabelle) is future work (Section 10).
% ===========================================================================
\documentclass[11pt]{article}
% ---------------------------------------------------------------------------
% preamble.tex -- packages, theorem environments, layout
% ---------------------------------------------------------------------------
\usepackage[T1]{fontenc}
\usepackage[utf8]{inputenc}
\usepackage{lmodern}
\usepackage[a4paper,margin=2.6cm]{geometry}
\usepackage{microtype}
\usepackage{amsmath,amssymb,amsthm,mathtools}
\usepackage{stmaryrd}
\usepackage{booktabs,tabularx,array}
\usepackage{enumitem}
\usepackage{xcolor}
\usepackage{listings}
\usepackage[hidelinks]{hyperref}
\usepackage[capitalise,noabbrev]{cleveref}

\setlist{itemsep=2pt,topsep=4pt}

\theoremstyle{definition}
\newtheorem{definition}{Definition}[section]
\newtheorem{assumption}[definition]{Assumption}
\newtheorem{example}[definition]{Example}
\theoremstyle{plain}
\newtheorem{theorem}[definition]{Theorem}
\newtheorem{lemma}[definition]{Lemma}
\newtheorem{proposition}[definition]{Proposition}
\newtheorem{corollary}[definition]{Corollary}
\theoremstyle{remark}
\newtheorem{remark}[definition]{Remark}

\crefname{assumption}{Assumption}{Assumptions}

\lstdefinestyle{kernel}{
  language=C++,
  basicstyle=\ttfamily\small,
  keywordstyle=\bfseries,
  commentstyle=\itshape\color{black!60},
  columns=fullflexible,
  frame=single,
  framerule=0.3pt,
  xleftmargin=4pt,
  xrightmargin=4pt,
  showstringspaces=false,
  escapeinside={(*}{*)},
}
\lstset{style=kernel}

% ---------------------------------------------------------------------------
% notation.tex -- macros used throughout the report
% ---------------------------------------------------------------------------
\newcommand{\Nat}{\mathbb{N}}
\newcommand{\Int}{\mathbb{Z}}
\newcommand{\Rat}{\mathbb{Q}}
\newcommand{\Real}{\mathbb{R}}
\newcommand{\Rnn}{\mathbb{R}_{\geq 0}}
% time domain and horizon
\newcommand{\TT}{\mathbb{T}}
\newcommand{\TR}{\mathbb{T}_R}
\newcommand{\Hz}{H}
% automata and models
\newcommand{\VD}{V_D}
\newcommand{\VP}{V_P}
\newcommand{\IR}{\mathit{IR}}
\newcommand{\Ksem}{K_{\mathit{sem}}}
\newcommand{\Kprod}{K_{\mathit{prod}}}
\newcommand{\Loc}{\mathit{L}}
\newcommand{\Clk}{\mathit{C}}
\newcommand{\Ch}{\mathit{Ch}}
\newcommand{\Act}{\mathit{Act}}
\newcommand{\Edges}{\mathit{E}}
\newcommand{\Inv}{\mathit{Inv}}
\newcommand{\Conf}{\mathit{Conf}}
\newcommand{\Valid}{\mathit{Conf}_{\checkmark}}
\newcommand{\Props}{\mathit{AP}}
\newcommand{\lab}{\mathcal{L}}
\newcommand{\at}[1]{\mathsf{at}(#1)}
\newcommand{\tauact}{\tau}
\newcommand{\cons}{\Phi}
\newcommand{\sem}[1]{\llbracket #1 \rrbracket}
% relations
\newcommand{\iso}{\cong}
\newcommand{\wbis}{\approx}
\newcommand{\sbis}{\sim}
\newcommand{\simul}{\preccurlyeq}
\newcommand{\trans}[1]{\xrightarrow{\,#1\,}}
\newcommand{\wtrans}[1]{\xRightarrow{\,#1\,}}
% kernel
\newcommand{\kappaconf}{\kappa}
\newcommand{\W}{\mathcal{W}}
\newcommand{\ticks}{\mathit{ticks}}
\newcommand{\Tmax}{T_{\max}}
\newcommand{\abs}{\alpha}
\newcommand{\code}[1]{\texttt{#1}}
\newcommand{\ok}{\mathsf{ok}}
\newcommand{\err}{\mathsf{err}}
\newcommand{\sat}{\models}


\title{A Verified Twin Execution Kernel:\\
Compiler and Kernel Correctness for Semantically Aligned Digital Twins}
\author{Technical report accompanying the Verified Twin repository}
\date{Report version 1.0 \quad (compiler 1.0.0, kernel 1.0.0, IR format \texttt{twin-ir/1})}

\begin{document}
\maketitle

\begin{abstract}
A Digital Twin (DT) view $\VD$ is a timed automaton that has been checked, by the
semantic aligner of \emph{SemPTDTAlignmentICSE}, to be semantically aligned with
the corresponding Physical Twin view $\VP$. This report proves that the Digital
Twin implementation $D$ produced by the Verified Twin toolchain realises $\VD$
\emph{by construction}. The implementation compiles $\VD$ into a canonical
intermediate representation $\IR(\VD)$, executes it in a minimal semantic kernel
$\Ksem$, and embeds the kernel in a production runtime $\Kprod$. We prove
\[
  \VD \;\iso\; \IR(\VD) \;\iso\; \Ksem(\IR(\VD)) \;\wbis\; \Kprod(\IR(\VD)),
  \qquad\text{hence}\qquad \VD \;\wbis\; \Kprod(\IR(\VD)),
\]
under an explicitly defined \emph{logical-time} semantics and explicitly stated
assumptions. The first two relations are isomorphisms of labelled timed
transition systems; the last is a weak timed bisimulation that abstracts
runtime bookkeeping (ledger writes, API serialisation, metrics, publication)
as internal $\tauact$-actions. We make no claim about hard real-time
behaviour, wall-clock time, schedulers or networks, and the proof has not
been mechanised.
\end{abstract}

\tableofcontents
\newpage

% ===========================================================================
\section{Introduction and preliminaries}\label{sec:intro}
% ===========================================================================

\paragraph{Context.}
The alignment paper~\cite{icse} establishes, at design time, the chain
\[
  P \;\simul\; \VP \;\sbis_{\cons}\; \VD \;\sbis\; D,
\]
where $P$ is the physical system, $\VP$ and $\VD$ are timed-automaton views of
the Physical and Digital Twin, $\sbis_{\cons}$ is semantic alignment under the
domain knowledge $\cons=(K,\{I_P,I_D\})$, and $D$ is the deployed Digital Twin.
The paper argues that the last relation, $\VD\sbis D$, can be obtained by
construction because the DT is software under the engineer's control. This
report makes that argument concrete for one implementation: the Verified Twin
toolchain of this repository.

\paragraph{The chain proved here.}
\begin{center}
\begin{tabular}{@{}lll@{}}
\toprule
Relation & Meaning & Result \\ \midrule
$\VD \iso \IR(\VD)$ & the compiler preserves and reflects the semantics & \cref{thm:compiler} \\
$\IR(\VD) \iso \Ksem(\IR(\VD))$ & the kernel executes exactly the IR semantics & \cref{thm:kernel} \\
$\Ksem \wbis \Kprod$ & runtime bookkeeping is unobservable ($\tauact$) & \cref{thm:production} \\
$\VD \wbis \Kprod(\IR(\VD))$ & composition & \cref{thm:composition} \\
\bottomrule
\end{tabular}
\end{center}

\paragraph{Method.}
The implementation and the proof were developed together. Wherever an
implementation choice would have complicated the argument, the
implementation was simplified instead. Examples are the exact integer time
grid (\cref{sec:time}), value-type kernel configurations that map one-to-one
onto formal configurations (\cref{sec:kernel}), and a compiler that
\emph{rejects}, rather than approximates, every construct whose semantics
it cannot preserve exactly (\cref{sec:compiler}). \Cref{app:code} maps every
definition to the C++ function that implements it.

\paragraph{What is not claimed.}
Time is \emph{logical model time}. No statement is made about wall-clock
time, worst-case execution time, scheduling, network latency, clock
synchronisation, the physical plant, or the correctness of planners and
sensor models (\cref{sec:limitations}). The argument is a pen-and-paper proof:
it has not been checked by a proof assistant. Tests in the repository are
validation, not proof.

% ---------------------------------------------------------------------------
\subsection{Labelled timed transition systems}
% ---------------------------------------------------------------------------

Throughout, $\TT\subseteq\Rnn$ denotes a \emph{time domain}: a set containing
$0$ and closed under addition, ordered by $\le$. The domain used by the
implementation is fixed in \cref{sec:time}.

\begin{definition}[LTTS]\label{def:ltts}
A \emph{labelled timed transition system} over $\TT$ is a tuple
$\mathcal{T}=(S,s_0,\Sigma,\to,\Props,\lab)$ where $S$ is a set of states,
$s_0\in S$ is the initial state, $\Sigma$ is a set of discrete labels with
$\Sigma\cap\TT=\emptyset$, ${\to}\subseteq S\times(\Sigma\cup\TT)\times S$ is the
transition relation, $\Props$ is a set of atomic propositions and
$\lab:S\to2^{\Props}$ is the labelling. A designated subset
$\Sigma_\tauact\subseteq\Sigma$ contains the \emph{internal} labels.
We write $s\trans{\lambda}s'$ for $(s,\lambda,s')\in{\to}$.
\end{definition}

This is Definition~1 of the alignment paper, with the event and delay labels
written $\Sigma$ and $\TT$.

\begin{definition}[Isomorphism]\label{def:iso}
Two LTTSs $\mathcal{T}_1,\mathcal{T}_2$ over the same $\TT$ are
\emph{isomorphic}, written $\mathcal{T}_1\iso\mathcal{T}_2$, if there are
bijections $\gamma:S_1\to S_2$, $\eta:\Sigma_1\to\Sigma_2$ (with
$\eta(\Sigma_{1,\tauact})=\Sigma_{2,\tauact}$) and $\pi:\Props_1\to\Props_2$ such that
(i) $\gamma(s_{0,1})=s_{0,2}$;
(ii) $s\trans{d}s'$ iff $\gamma(s)\trans{d}\gamma(s')$ for all $d\in\TT$;
(iii) $s\trans{a}s'$ iff $\gamma(s)\trans{\eta(a)}\gamma(s')$ for all $a\in\Sigma_1$;
(iv) $p\in\lab_1(s)$ iff $\pi(p)\in\lab_2(\gamma(s))$.
\end{definition}

\begin{definition}[Weak timed bisimulation]\label{def:wtb}
Let $\mathcal{T}_1,\mathcal{T}_2$ be LTTSs over $\TT$ whose observable label
sets are identified by a bijection $\eta$ and whose propositions are
identified by $\pi$. For $\mathcal{T}_i$ write
$s\wtrans{\epsilon}s'$ if $s$ reaches $s'$ by a finite sequence of internal
transitions, $s\wtrans{a}s'$ if $s\wtrans{\epsilon}\trans{a}\wtrans{\epsilon}s'$
for an observable $a$, and $s\wtrans{d}s'$ if
$s\wtrans{\epsilon}\trans{d_1}\wtrans{\epsilon}\cdots\trans{d_k}\wtrans{\epsilon}s'$
with $d_1+\dots+d_k=d$. A relation $\mathcal{R}\subseteq S_1\times S_2$ is a
\emph{weak timed bisimulation} if $(s_{0,1},s_{0,2})\in\mathcal{R}$ and for
every $(s_1,s_2)\in\mathcal{R}$:
\begin{enumerate}[label=(W\arabic*)]
\item $p\in\lab_1(s_1)$ iff $\pi(p)\in\lab_2(s_2)$;
\item if $s_1\trans{\lambda}s_1'$ ($\lambda$ observable or a delay) there is
  $s_2\wtrans{\eta(\lambda)}s_2'$ with $(s_1',s_2')\in\mathcal{R}$, and if
  $s_1\trans{\lambda}s_1'$ with $\lambda$ internal there is
  $s_2\wtrans{\epsilon}s_2'$ with $(s_1',s_2')\in\mathcal{R}$;
\item symmetrically for transitions of $s_2$.
\end{enumerate}
$\mathcal{T}_1\wbis\mathcal{T}_2$ if such a relation exists. Taking
$\wtrans{}$ to be $\trans{}$ yields \emph{strong} timed bisimulation $\sbis$.
\end{definition}

\begin{lemma}[Standard facts]\label{lem:facts}
(a) $\mathcal{T}_1\iso\mathcal{T}_2$ implies $\mathcal{T}_1\sbis\mathcal{T}_2$
(the graph of $\gamma$ is a strong bisimulation), and strong timed bisimilarity
implies weak timed bisimilarity.
(b) $\iso$, $\sbis$ and $\wbis$ are equivalence relations; in particular they
are transitive (relational composition of bisimulations is a bisimulation).
(c) Isomorphic LTTSs satisfy the same properties expressible over their
labels and propositions.
\end{lemma}
\begin{proof}
(a) Conditions (i)--(iv) of \cref{def:iso} are exactly the strong-bisimulation
conditions for the relation $\{(s,\gamma(s))\}$, and $\trans{}\subseteq\wtrans{}$.
(b) is standard~\cite{milner}; for weak timed bisimulation the composition
argument also covers delays because $\wtrans{d}$ composes additively.
(c) follows because $\gamma$ maps runs to runs bijectively and preserves labels
and propositions.
\end{proof}

% ===========================================================================
\section{Source models and their logical-time semantics}\label{sec:source}
% ===========================================================================

\subsection{The supported fragment}

The source $\VD$ is a UPPAAL XML document. Its meaning is the abstract syntax
that the UPPAAL parser UTAP builds from it, which is the same parser the semantic
aligner uses. The compiler accepts exactly the fragment the aligner gives
semantics to (see \texttt{docs/existing-aligner-integration.md}); everything
else is rejected (\cref{sec:compiler}).

\begin{definition}[Clock constraints]\label{def:constraints}
Let $\Clk$ be a finite set of clocks and $0\notin\Clk$ a \emph{reference clock}.
An \emph{atomic constraint} is $x-y\bowtie c$ with $x\in\Clk$,
$y\in\Clk\cup\{0\}$, $x\neq y$, ${\bowtie}\in\{<,\le,=,\ge,>\}$ and $c\in\Int$;
for $y=0$ it is written $x\bowtie c$ (a \emph{simple} constraint), otherwise it is
a \emph{diagonal} constraint. A \emph{conjunction} is a finite \emph{set} $G$ of
atomic constraints. A valuation $v:\Clk\to\TT$ is extended by $v(0)=0$, and
\[
  v\sat(x-y\bowtie c)\iff v(x)-v(y)\bowtie c,\qquad
  v\sat G\iff \forall\varphi\in G.\;v\sat\varphi .
\]
$\cons(\Clk)$ denotes all conjunctions, $\cons_s(\Clk)$ those with simple atoms only.
\end{definition}

\begin{definition}[Source automaton]\label{def:source}
A \emph{source automaton} is $A=(\Loc,\ell_0,\Clk,\Ch,\Edges,\Inv)$ where
$\Loc=(\ell_0',\dots,\ell_{k-1}')$ is a finite sequence of distinctly named
locations (document order) with $\ell_0\in\Loc$; $\Clk=(x_1,\dots,x_n)$ is a
finite sequence of distinct clocks (global declarations first, then template
declarations, both in declaration order); $\Ch$ is a finite set of channels;
$\Act=\{\tauact\}\cup\{a!,a?\mid a\in\Ch\}$;
$\Edges=(e_1,\dots,e_m)$ is a finite sequence of edges (document order)
$e_j=(\ell,g,\alpha,Y,\ell')$ with $\ell,\ell'\in\Loc$, $g\in\cons_s(\Clk)$,
$\alpha\in\Act$, $Y\subseteq\Clk$; and $\Inv:\Loc\to\cons(\Clk)$.
\end{definition}

Edges are identified by their position, so two edges with equal components are
still distinct transitions (UPPAAL semantics). Guards are simple conjunctions
because the aligner's guard reader cannot read diagonals (finding~6 of the
integration report), whereas invariants may contain diagonals.

\begin{remark}[Side conditions enforced by the compiler]\label{rem:side}
In addition, the compiler requires: exactly one template, instantiated once;
no data variables, functions, select bindings, urgent/committed locations,
urgent/broadcast channels, priorities or probabilistic weights; resets to
$0$ only; constants within UDBM's bound range; and $\ell_0=\ell_0'$, that is,
the initial location is the first location in document order, which is the
aligner's assumption (finding~4). Under these conditions the abstract syntax
read by the compiler and the one read by the aligner coincide (\cref{prop:tv}).
\end{remark}

\subsection{Logical time}

\begin{definition}[Logical-time grid and horizon]\label{def:grid}
Fix a resolution $R\in\{10^0,10^1,\dots,10^9\}$ (ticks per model time unit) and
the horizon $\Hz=\Tmax/R$ with $\Tmax=2^{62}-1$. The \emph{logical time domain} is
\[
  \TR=\{\,k/R\mid k\in\Nat\,\},\qquad \TR^{\Hz}=\{\,t\in\TR\mid t\le\Hz\,\}.
\]
\end{definition}

$\TR$ contains $0$ and is closed under addition, so it is a time domain in the
sense of \cref{sec:intro}. It consists exactly of the non-negative decimal
numbers with at most $\log_{10}R$ fractional digits, which makes timestamps such
as $31.5$ representable exactly. Values off the grid are rejected, never rounded
(\cref{sec:time}). At $R=1000$ the horizon $\Hz$ is about $4.6\cdot10^{15}$ time
units.

\subsection{Semantics}

\begin{definition}[Configurations]\label{def:conf}
A \emph{configuration} of $A$ is a triple $(\ell,v,\theta)$ with $\ell\in\Loc$,
$v:\Clk\to\TR$ and $\theta\in\TR^{\Hz}$ such that $v(x)\le\theta$ for all $x$.
$\theta$ is the \emph{global logical time}, an implicit clock that is never
reset. The set of configurations is $\Conf(A)$; a configuration is
\emph{valid} if $v\sat\Inv(\ell)$; the valid configurations form $\Valid(A)$.
For $d\in\TR$, $(v+d)(x)=v(x)+d$; for $Y\subseteq\Clk$, $v[Y{:=}0](x)=0$ if $x\in Y$
and $v(x)$ otherwise.
\end{definition}

\begin{definition}[Logical-time semantics $\sem{A}$]\label{def:source-sem}
The LTTS $\sem{A}=(\Valid(A),s_0,\Edges,\to,\Props_A,\lab_A)$ over $\TR$ has
internal labels $\{e\in\Edges\mid \mathit{act}(e)=\tauact\}$,
$\Props_A=\{\at{\ell}\mid\ell\in\Loc\}$ and $\lab_A(\ell,v,\theta)=\{\at{\ell}\}$.
The initial state and the transitions are given by the rules
\begin{align*}
&\textsc{(Init)} && s_0=(\ell_0,\mathbf{0},0) \quad\text{provided } \mathbf{0}\sat\Inv(\ell_0);\\
&\textsc{(Delay)} && (\ell,v,\theta)\trans{d}(\ell,v+d,\theta+d)
  \quad\text{iff } \theta+d\le\Hz \text{ and } \forall d'\in[0,d]\subseteq\Rnn.\;v+d'\sat\Inv(\ell);\\
&\textsc{(Disc)} && (\ell,v,\theta)\trans{e}(\ell',v[Y{:=}0],\theta)
  \quad\text{iff } e=(\ell,g,\alpha,Y,\ell'),\; v\sat g,\; v[Y{:=}0]\sat\Inv(\ell').
\end{align*}
$A$ is \emph{initialisable} if $\mathbf{0}\sat\Inv(\ell_0)$; all results below
are about initialisable automata (the kernel refuses to start otherwise).
The \emph{action-labelled} view of $\sem{A}$ is obtained by relabelling each
edge $e$ with $\mathit{act}(e)$; it is the timed transition system of the
alignment paper, with $\tauact$-edges internal.
\end{definition}

The intermediate points $d'$ in \textsc{(Delay)} range over the \emph{reals}: the
invariant must hold continuously, as in the dense-time semantics of timed
automata~\cite{alur-dill,uppaal}. The next lemma shows that this condition can be
checked at the two end points.

\begin{lemma}[Convexity of delays]\label{lem:convex}
For every conjunction $G$ and valuation $v$, the set
$D(v,G)=\{d'\in\Rnn\mid v+d'\sat G\}$ is an interval. Consequently
\[
  \forall d'\in[0,d].\;v+d'\sat G \iff v\sat G\ \wedge\ v+d\sat G .
\]
\end{lemma}
\begin{proof}
For a diagonal atom $x-y\bowtie c$, $(v(x)+d')-(v(y)+d')=v(x)-v(y)$ does not depend
on $d'$, so its solution set is $\emptyset$ or $\Rnn$. For a simple atom
$x\bowtie c$ the solution set $\{d'\mid v(x)+d'\bowtie c\}$ is a (possibly empty
or degenerate) half-line or point of $\Rnn$. Hence $D(v,G)$ is a finite
intersection of intervals, an interval. If an interval contains $0$ and $d$ it
contains $[0,d]$; the converse direction is trivial.
\end{proof}

\begin{lemma}[Well-formedness is invariant]\label{lem:wf}
Every state reachable from $s_0$ in $\sem{A}$ is a valid configuration
$(\ell,v,\theta)$ with $0\le v(x)\le\theta\le\Hz$ for all $x\in\Clk$.
\end{lemma}
\begin{proof}
By induction on run length. $s_0$ satisfies it by \textsc{(Init)}. \textsc{(Delay)}
adds $d$ to every clock and to $\theta$, preserving $v(x)\le\theta$, enforces
$\theta+d\le\Hz$, and requires $v+d\sat\Inv(\ell)$. \textsc{(Disc)} keeps $\theta$, sets
some clocks to $0\le\theta$, and requires the target invariant.
\end{proof}

\begin{proposition}[Relation to dense time]\label{prop:dense}
Let $\sem{A}_{\Real}$ be the standard dense-time semantics: the same rules over
$\TT=\Rnn$ with no horizon. Then $\sem{A}$ is the sub-LTTS of $\sem{A}_{\Real}$
induced by the configurations with values in $\TR$ and $\theta\le\Hz$, and by
delays in $\TR$. That is, for grid configurations $s,s'$ and $\lambda\in\Edges\cup\TR$:
$s\trans{\lambda}s'$ in $\sem{A}$ iff $s\trans{\lambda}s'$ in $\sem{A}_{\Real}$ and
$\theta(s')\le\Hz$. In particular every timed run of $\sem{A}$ is a timed run of
$\sem{A}_{\Real}$.
\end{proposition}
\begin{proof}
The rules are literally the same; \textsc{(Delay)} quantifies over real
intermediate points in both semantics. A grid configuration plus a grid delay
yields a grid configuration, and resets yield grid values.
\end{proof}

\begin{remark}
\Cref{prop:dense} says the grid semantics \emph{under-approximates} dense time
soundly. The grid cannot reproduce dense-time behaviours that require
timestamps off the grid, for example two events $10^{-4}$ time units apart at
$R=1000$. Such behaviours are not representable in the logical time of the
implementation, by design. For \emph{closed} automata (no strict constraints)
integer time already preserves location reachability~\cite{hmp}. For automata
with strict constraints, a finer grid may be needed. We do not claim a general
completeness bound for the grid; this remark is informative and is not used in
the proofs.
\end{remark}

% ===========================================================================
\section{The Twin Intermediate Representation}\label{sec:ir}
% ===========================================================================

The IR (\code{include/twin/ir/model.hpp}) is an executable canonicalisation of a
source automaton. It is deliberately not a programming language: it contains
exactly the information of \cref{def:source}, with names replaced by dense
indices, conjunctions in canonical order, and explicit propositions.

\begin{definition}[IR model]\label{def:ir}
An IR model is a tuple
$M=(\mathit{info},R,X,\Ch_M,\Lambda,\iota,T,P)$ where
\begin{itemize}
\item $\mathit{info}$ holds identity and provenance (model id, version, source
  template, SHA-256 of the source); it has no semantic content;
\item $R$ is the logical-time resolution of \cref{def:grid};
\item $X=(X_1,\dots,X_n)$ are clock names; clock \emph{identifier} $i\in\{1..n\}$
  denotes $X_i$, and identifier $0$ the reference clock;
\item $\Ch_M$ is a sorted list of channel names;
\item $\Lambda=(\lambda_0,\dots,\lambda_{k-1})$ are locations
  $\lambda_i=(\mathit{id}_i,J_i)$ with a canonical conjunction $J_i$ (the invariant);
\item $\iota\in\{0,\dots,k-1\}$ is the initial location;
\item $T=(t_0,\dots,t_{m-1})$ are transitions
  $t_j=(\mathit{tid}_j,\mathit{src}_j,\mathit{tgt}_j,\alpha_j,g_j,Y_j)$ with location indices,
  an action $\alpha_j\in\{\tauact\}\cup\{a!,a?\mid a\in\Ch_M\}$, a canonical conjunction
  $g_j$ and a sorted, duplicate-free reset list $Y_j\subseteq\{1..n\}$;
\item $P=(p_0,\dots,p_{k-1})$ are propositions $p_i=(\at{\mathit{id}_i},i,\phi_i)$ where
  $\phi_i$ is the ontology formula $I_D(\mathit{id}_i)$ (possibly empty). The formula is
  metadata: it does not influence the transition relation.
\end{itemize}
Atomic constraints are quadruples $(i,j,\bowtie,c)$ with $1\le i\le n$, $0\le j\le n$,
$i\ne j$, meaning $X_i-X_j\bowtie c$ (simple if $j=0$). A conjunction is
\emph{canonical} if it is sorted strictly increasingly in the lexicographic order
on quadruples. That makes it a set with a unique list representation.
\end{definition}

\begin{definition}[Well-formed IR]\label{def:wf-ir}
$M$ is well-formed if identifiers are syntactically valid and unique
(locations, clocks, channels, transitions); all indices are in range;
conjunctions are canonical and their constants lie in $[-(2^{31}-1),2^{31}-1]$;
reset lists are sorted, unique and do not contain $0$; every non-internal action
uses a declared channel; $|P|=k$ with $p_i=(\at{\mathit{id}_i},i,\cdot)$; and $R$ is
a power of ten in $[1,10^9]$. These rules are implemented by
\code{twin::ir::validate} (\code{src/ir/validate.cpp}).
\end{definition}

\begin{definition}[IR semantics $\sem{M}$]\label{def:ir-sem}
Configurations of $M$ are triples $(i,w,\theta)$ with $i\in\{0..k-1\}$,
$w:\{1..n\}\to\TR$ and $\theta\in\TR^{\Hz}$ with $w(x)\le\theta$. Satisfaction of
$(i,j,\bowtie,c)$ by $w$ is $w(i)-w(j)\bowtie c$ with $w(0)=0$. The LTTS $\sem{M}$ over
$\TR$ has as states the valid configurations ($w\sat J_i$), the labels $T$, internal
labels $\{t_j\mid\alpha_j=\tauact\}$, propositions $P$ with
$\lab_M(i,w,\theta)=\{p_i\}$, and the rules
\begin{align*}
&\textsc{(Init$_M$)} && s_0=(\iota,\mathbf{0},0)\quad\text{provided }\mathbf{0}\sat J_\iota;\\
&\textsc{(Delay$_M$)} && (i,w,\theta)\trans{d}(i,w+d,\theta+d)
  \quad\text{iff } \theta+d\le\Hz,\ w\sat J_i,\ w+d\sat J_i;\\
&\textsc{(Disc$_M$)} && (i,w,\theta)\trans{t_j}(\mathit{tgt}_j,w[Y_j{:=}0],\theta)
  \quad\text{iff } i=\mathit{src}_j,\ w\sat g_j,\ w[Y_j{:=}0]\sat J_{\mathit{tgt}_j}.
\end{align*}
\end{definition}

By \cref{lem:convex}, \textsc{(Delay$_M$)} is equivalent to the continuous
formulation of \textsc{(Delay)}. The IR rule is written with end points because
that is the form the kernel implements.

\subsection{Canonical serialisation}

\begin{definition}[Canonical text]\label{def:canon}
$\mathit{canon}(M)$ is the JSON document of format \code{twin-ir/1}
(\code{include/twin/ir/codec.hpp}) serialised with object keys sorted bytewise,
no insignificant whitespace, and integers only (no floating-point values). The
IR hash is $h(M)=\mathrm{SHA\text{-}256}(\mathit{canon}(M))$~\cite{fips}. For the
documents used here this coincides with RFC~8785~\cite{jcs}.
\end{definition}

\begin{lemma}[Canonicity and faithful loading]\label{lem:canon}
For well-formed $M,M'$: $\mathit{canon}(M)=\mathit{canon}(M')$ iff $M=M'$. Moreover,
the loader \code{from\_canonical\_text} accepts a byte string $b$ only if it decodes
$b$ into a well-formed $M$ with $\mathit{canon}(M)=b$.
\end{lemma}
\begin{proof}
($\Leftarrow$) Serialisation is a function. ($\Rightarrow$) Every component of $M$
is written explicitly; clocks, channels, locations, transitions and propositions
are emitted as arrays in model order; references use unique identifiers; and
conjunction and reset lists are canonical, so their list order is determined by
the set. Hence $M$ can be recovered from $\mathit{canon}(M)$, and canon is
injective. The second claim is the round-trip check implemented in
\code{from\_canonical\_text}: the decoded model is validated, re-serialised and
compared byte for byte with the input.
\end{proof}

\begin{remark}[Why the round trip matters]
\Cref{lem:canon} reduces trust in the JSON parser. Whatever the parser does,
the kernel only ever executes a model whose canonical serialisation is exactly
the hashed byte string recorded in the package manifest. Duplicate keys,
alternative number spellings and whitespace variants are rejected.
\end{remark}

% ===========================================================================
\section{Compiler correctness: $\VD\iso\IR(\VD)$}\label{sec:compiler}
% ===========================================================================

\subsection{The translation}

\begin{definition}[Compiler]\label{def:compiler}
Let $A=(\Loc,\ell_0,\Clk,\Ch,\Edges,\Inv)$ satisfy \cref{rem:side}, with
$\Loc=(\ell_0',\dots,\ell_{k-1}')$, $\Clk=(x_1,\dots,x_n)$ and $\Edges=(e_1,\dots,e_m)$. Define
the index maps
\[
  \nu_\Loc(\ell_i')=i,\qquad \nu_\Clk(x_i)=i,\ \nu_\Clk(0)=0,\qquad \nu_\Edges(e_j)=t_{j-1},
\]
and lift $\nu_\Clk$ to atoms and conjunctions:
$\nu(x-y\bowtie c)=(\nu_\Clk(x),\nu_\Clk(y),\bowtie,c)$ and $\nu(G)=\{\nu(\varphi)\mid\varphi\in G\}$.
Let $\mathit{sort}$ denote the canonical list of a set. Then $C(A)=M$ with
\begin{itemize}
\item $X=(x_1,\dots,x_n)$, $\Ch_M=\mathit{sort}(\Ch)$, $R$ and $\mathit{info}$ from the options;
\item $\lambda_i=(\mathrm{name}(\ell_i'),\ \mathit{sort}(\nu(\Inv(\ell_i'))))$;
\item $\iota=\nu_\Loc(\ell_0)$ $(=0$ by \cref{rem:side});
\item for $e_j=(\ell,g,\alpha,Y,\ell')$:
  $t_{j-1}=(\mathit{tid}_{j-1},\nu_\Loc(\ell),\nu_\Loc(\ell'),\alpha,\mathit{sort}(\nu(g)),\mathit{sort}(\nu_\Clk(Y)))$,
  where $\mathit{tid}$ is ``$\mathit{src}.\alpha.\mathit{tgt}$'', suffixed by $\#r$ for the
  $r$-th repetition of the same triple in document order;
\item $p_i=(\at{\mathrm{name}(\ell_i')},i,I_D(\ell_i'))$, with $I_D(\ell)$ the entry of the
  DT interpretation file or the empty string.
\end{itemize}
This is \code{twin::compiler::compile\_file} (\code{src/compiler/compiler.cpp}).
\end{definition}

\begin{lemma}\label{lem:nu}
$\nu_\Loc$, $\nu_\Clk$ and $\nu_\Edges$ are bijections onto the locations, clock
identifiers and transitions of $C(A)$, and $C(A)$ is well-formed.
\end{lemma}
\begin{proof}
One IR location, clock and transition is emitted per source location, clock
and edge, in the same order. Location and clock names are unique in $A$
(UTAP rejects duplicates, and the compiler re-checks), so the inverse maps are
well defined. Transition identifiers are unique because repetitions are
numbered. Well-formedness is checked by \code{ir::validate} before the IR is
emitted. A failure there would be a compiler defect, reported as diagnostic
\code{TWC090}, and no IR would be produced.
\end{proof}

\begin{lemma}[Constraint preservation]\label{lem:constraint}
For every conjunction $G$ over $\Clk$ and valuation $v$:
$v\sat G \iff (v\circ\nu_\Clk^{-1})\sat\mathit{sort}(\nu(G))$.
\end{lemma}
\begin{proof}
Let $w=v\circ\nu_\Clk^{-1}$, so $w(\nu_\Clk(x))=v(x)$, and $w(0)=v(0)=0$. For an atom,
$w(\nu_\Clk(x))-w(\nu_\Clk(y))=v(x)-v(y)$, so $v\sat\varphi\iff w\sat\nu(\varphi)$.
Satisfaction of a conjunction depends only on the \emph{set} of its atoms, and
$\mathit{sort}$ does not change that set.
\end{proof}

\subsection{The isomorphism}

\begin{theorem}[Compiler correctness]\label{thm:compiler}
Let $A$ satisfy \cref{rem:side}, be initialisable, and let $M=C(A)$. The map
\[
  \Gamma:\Valid(A)\to\Valid(M),\qquad
  \Gamma(\ell,v,\theta)=(\nu_\Loc(\ell),\ v\circ\nu_\Clk^{-1},\ \theta)
\]
is a bijection, and $(\Gamma,\nu_\Edges,\pi)$ with $\pi(\at{\ell})=p_{\nu_\Loc(\ell)}$ is an
isomorphism $\sem{A}\iso\sem{M}$ (\cref{def:iso}). The isomorphism maps internal
labels to internal labels and preserves action labels:
$\mathit{act}(e)=\alpha(\nu_\Edges(e))$.
\end{theorem}
\begin{proof}
\emph{Bijection.} $\Gamma$ is the product of bijections (\cref{lem:nu}) with
identity on $\theta$. The side conditions $v(x)\le\theta\le\Hz$ are
invariant under renaming. Validity is preserved and reflected by
\cref{lem:constraint} applied to $\Inv(\ell)$ and $J_{\nu_\Loc(\ell)}=\mathit{sort}(\nu(\Inv(\ell)))$.

\emph{(i) Initial state.} $\Gamma(\ell_0,\mathbf{0},0)=(\nu_\Loc(\ell_0),\mathbf{0},0)=(\iota,\mathbf{0},0)$.
$A$ is initialisable iff $M$ is, by \cref{lem:constraint}.

\emph{(ii) Delays.} By \cref{lem:convex}, $(\ell,v,\theta)\trans{d}(\ell,v+d,\theta+d)$
holds iff $\theta+d\le\Hz$, $v\sat\Inv(\ell)$ and $v+d\sat\Inv(\ell)$. By
\cref{lem:constraint}, and since $(v+d)\circ\nu_\Clk^{-1}=(v\circ\nu_\Clk^{-1})+d$, this
is exactly the premise of \textsc{(Delay$_M$)} for $\Gamma(\ell,v,\theta)$. The targets
correspond: $\Gamma(\ell,v+d,\theta+d)=(\nu_\Loc(\ell),(v\circ\nu_\Clk^{-1})+d,\theta+d)$.

\emph{(iii) Discrete steps.} Let $e=(\ell,g,\alpha,Y,\ell')$ and $t=\nu_\Edges(e)$.
The premises of \textsc{(Disc)} are $v\sat g$ and $v[Y{:=}0]\sat\Inv(\ell')$. Those of
\textsc{(Disc$_M$)} at $\Gamma(\ell,v,\theta)$ are $\nu_\Loc(\ell)=\mathit{src}(t)$, which holds
by construction, $w\sat\mathit{sort}(\nu(g))$, and $w[\nu_\Clk(Y){:=}0]\sat J_{\nu_\Loc(\ell')}$.
They coincide by \cref{lem:constraint}, using
$v[Y{:=}0]\circ\nu_\Clk^{-1}=(v\circ\nu_\Clk^{-1})[\nu_\Clk(Y){:=}0]$. The targets
correspond in the same way. Conversely, every $t\in T$ is $\nu_\Edges(e)$ for a
unique $e$, so every IR step is the image of a source step.

\emph{(iv) Propositions.} $\lab_A(\ell,v,\theta)=\{\at{\ell}\}$ and
$\lab_M(\Gamma(\ell,v,\theta))=\{p_{\nu_\Loc(\ell)}\}=\{\pi(\at{\ell})\}$.

Actions are copied verbatim, so internal edges map to internal transitions.
\end{proof}

\subsection{From the abstract compiler to the implemented one}

\Cref{thm:compiler} is about the mathematical function $C$. Two links connect
it to the bytes on disk.

\begin{assumption}[Source reading]\label{ass:utap}
UTAP parses the UPPAAL document into the abstract syntax of
\cref{def:source}. This is the same assumption the alignment verdict rests on,
because the aligner uses the same parser.
\end{assumption}

\begin{proposition}[Translation validation]\label{prop:tv}
Suppose compilation succeeds. Then the automaton $A_c$ read by the compiler and
the automaton $A_a$ read by the aligner's \code{dtpta::TimedAutomaton} agree on:
the template, the clock sequence and numbering, the location names and their
order, the initial location (index $0$), the edge sequence (endpoints, actions,
reset sets), and, for every location and edge, the \emph{set of valuations}
satisfying the invariant and the guard. Consequently $\sem{A_c}=\sem{A_a}$.
\end{proposition}
\begin{proof}
These are exactly the checks of \code{crosscheck\_with\_aligner}
(\code{src/compiler/crosscheck.cpp}). Compilation fails with \code{TWC060} if any
of them fails. Invariants and guards are compared as zones: both readings are
intersected with the universe of non-negative valuations, closed, and compared
with UDBM's \code{dbm\_areEqual}. Two closed, non-empty DBMs are equal iff they
denote the same zone; empty results are compared as empty. The semantics of
\cref{def:source-sem} depends on conjunctions only through their satisfying
valuations, and all other components are equal, so the two LTTSs coincide.
\end{proof}

\begin{corollary}\label{cor:compiled-is-verified}
If compilation succeeds, then $\sem{\IR(\VD)}\iso\sem{A_a}$, where $A_a$ is the
automaton the aligner analysed (up to the aligner's own zone semantics; see
\cref{sec:limitations}).
\end{corollary}
\begin{proof}
\Cref{thm:compiler} gives $\sem{A_c}\iso\sem{M}$ and \cref{prop:tv} gives
$\sem{A_c}=\sem{A_a}$.
\end{proof}

\begin{remark}[Determinism of compilation]
$C$ is a function of the source bytes, the interpretation file and the options
$(\text{model id},\text{version},R)$. The implementation uses no timestamps,
paths, environment variables or iteration over unordered containers in the IR
(the compilation \emph{manifest} contains a timestamp and is not hashed into the
IR). Recompiling therefore yields a byte-identical IR. The test suite checks
this on the aligner's corpus.
\end{remark}

% ===========================================================================
\section{The semantic kernel $\Ksem$}\label{sec:kernel}
% ===========================================================================

The kernel (\code{include/twin/kernel}, \code{src/kernel}) is the only component
that computes semantic successors. Its library target links only the IR value
types and the dependency-free core (errors, exact time). It performs no I/O,
networking, persistence or logging. This section defines the kernel transition
system as the functions the code computes.

\subsection{Executable model}

\begin{definition}[Kernel model]\label{def:kmodel}
\code{Model::create(M)} succeeds for a well-formed $M$ (it re-runs
\code{ir::validate}) and produces $\hat M$, which consists of $M$ unchanged plus, for
every atom $(i,j,\bowtie,c)$ of every invariant and guard, the \emph{tick atom}
$(i,j,\bowtie,\hat c)$ with $\hat c=c\cdot R$. The multiplication is checked: if
$|c|>2^{31}-1$ or the product overflows, creation fails. The remaining data
(outgoing-transition lists sorted by index, a label table) are derived indexes
that change no semantics.
\end{definition}

\begin{definition}[Kernel configurations]\label{def:kconf}
A kernel configuration is the C++ value
\begin{lstlisting}
struct Configuration {
    uint32_t             location;  // l
    std::vector<int64_t> clocks;    // v, in ticks; clocks[i-1] is clock i
    int64_t              time;      // theta, in ticks
};
\end{lstlisting}
Mathematically, $\kappa=(l,c,\tau)\in\Nat\times\Int^n\times\Int$. Its value of clock
$i$ is $c_{i-1}$ for $i\ge1$ and $0$ for the reference clock. The set of
\emph{admissible} kernel configurations is
\[
  \W(\hat M)=\{(l,c,\tau)\mid l<k,\ |c|=n,\ 0\le c_{i}\le\tau\le\Tmax,\ (l,c,\tau)\text{ satisfies }\hat J_l\}.
\]
This is exactly what \code{check\_configuration} accepts.
\end{definition}

\subsection{Kernel functions}

The following functions are those of \code{src/kernel/semantics.cpp}. They are
pure: they read their arguments and the immutable $\hat M$ and return a value
or an error, and they have no other effect.

\begin{definition}[Atom satisfaction]\label{def:katom}
$\kappa$ satisfies the tick atom $(i,j,\bowtie,\hat c)$ iff
$\mathit{val}_\kappa(i)-\mathit{val}_\kappa(j)\bowtie\hat c$, computed in \code{int64}
(\code{atom\_holds}). $\kappa$ satisfies a conjunction iff it satisfies every atom
(\code{satisfies}).
\end{definition}

\begin{definition}[Kernel transition functions]\label{def:kfun}\leavevmode
\begin{lstlisting}
initial_configuration(M) =
  c0 := (iota, 0^n, 0);   ok(c0) if c0 satisfies J_iota else err
delay(M, (l,c,t), d) =
  if d < 0                              -> err(InvalidArgument)
  if (l,c,t) does not satisfy J_l        -> err(InvariantViolation)
  t' := t + d; c' := c + d   (checked; any value > T_max -> err(Overflow))
  if (l,c',t') does not satisfy J_l      -> err(InvariantViolation)
  ok((l,c',t'))
fire(M, (l,c,t), j) =                    // transition t_j = (id,src,tgt,a,g,Y)
  if j >= m                              -> err(InvalidArgument)
  if l != src                            -> err(NotEnabled)
  if (l,c,t) does not satisfy g          -> err(NotEnabled)
  c' := c with c'[y-1] := 0 for y in Y
  if (tgt,c',t) does not satisfy J_tgt   -> err(NotEnabled)
  ok((tgt,c',t))
propositions(M, (l,c,t)) = { i | p_i.location == l }      // = {l}
\end{lstlisting}
\end{definition}

\begin{definition}[Kernel LTTS $\Ksem(M)$]\label{def:ksem}
$\Ksem(M)$ is the LTTS over the tick domain with states $\W(\hat M)$, initial state
the result of \code{initial\_configuration}, labels $\{0,\dots,m-1\}$ (internal
iff $\alpha_j=\tauact$), propositions $P$, labelling \code{propositions}, and
\[
  \kappa\trans{\delta}\kappa' \iff \code{delay}(\hat M,\kappa,\delta)=\ok(\kappa'),
  \qquad
  \kappa\trans{j}\kappa' \iff \code{fire}(\hat M,\kappa,j)=\ok(\kappa').
\]
The delay label $\delta\in\Nat$ is in ticks. When kernel and IR are compared,
$\delta$ corresponds to the delay $\delta/R\in\TR$ (\cref{sec:iso}).
\end{definition}

\begin{remark}[Derived kernel functions]
\code{enabled\_transitions}, \code{discrete\_successors}, \code{timed\_step}
($=\code{fire}\circ\code{delay}$), \code{enabling\_window}, \code{observe} and the
exploration functions are defined in terms of \code{delay} and \code{fire} or
characterised by them (\cref{lem:window,prop:monitor}). They add no transitions.
\end{remark}

\subsection{Arithmetic safety}

\begin{lemma}[No overflow]\label{lem:overflow}
For $\kappa\in\W(\hat M)$, every \code{int64} expression evaluated by
\code{atom\_holds}, \code{delay} and \code{fire} is computed without overflow,
or the function returns the explicit error \code{ArithmeticOverflow}.
\end{lemma}
\begin{proof}
Clock values lie in $[0,\Tmax]$ with $\Tmax<2^{62}$, so differences lie in
$(-2^{62},2^{62})$. Tick bounds satisfy $|\hat c|\le(2^{31}-1)\cdot10^9<2^{61}$.
Comparisons of such values are exact. In \code{delay}, the additions $\tau+\delta$
and $c_i+\delta$ use \code{\_\_builtin\_add\_overflow}, and results above $\Tmax$ are
rejected explicitly. \code{fire} assigns only the constant $0$.
\end{proof}

% ===========================================================================
\section{Kernel correctness: $\IR(\VD)\iso\Ksem(\IR(\VD))$}\label{sec:iso}
% ===========================================================================

\begin{definition}[Abstraction function]\label{def:alpha}
\[
  \abs:\W(\hat M)\to\Valid(M),\qquad
  \abs(l,c,\tau)=\bigl(l,\ i\mapsto c_{i-1}/R,\ \tau/R\bigr).
\]
\end{definition}

\begin{lemma}[$\abs$ is a bijection]\label{lem:alpha}
$\abs$ is well defined and bijective. Its inverse is
$\abs^{-1}(i,w,\theta)=(i,(R\,w(1),\dots,R\,w(n)),R\,\theta)$.
\end{lemma}
\begin{proof}
Grid values are exactly the numbers $k/R$ with $k\in\Nat$, so division by $R$ is a
bijection between $\{0,\dots,\Tmax\}$ and $\TR^{\Hz}$. The constraints
$0\le c_i\le\tau\le\Tmax$ correspond to $0\le w(i)\le\theta\le\Hz$. The location index is
copied. Validity corresponds by \cref{lem:atom} below.
\end{proof}

\begin{lemma}[Atom correspondence]\label{lem:atom}
For $\kappa\in\W(\hat M)$ and every atom $(i,j,\bowtie,c)$ of $M$ with tick atom
$(i,j,\bowtie,cR)$:
\[
  \kappa \text{ satisfies } (i,j,\bowtie,cR) \iff \abs(\kappa)\sat(i,j,\bowtie,c).
\]
\end{lemma}
\begin{proof}
Write $\abs(\kappa)=(l,w,\theta)$, so $w(i)=\mathit{val}_\kappa(i)/R$ for all $i$, including the
reference clock with value $0$. Then $\mathit{val}_\kappa(i)-\mathit{val}_\kappa(j)\bowtie cR$ iff
$R\,(w(i)-w(j))\bowtie R\,c$ iff $w(i)-w(j)\bowtie c$, because multiplication by $R>0$
preserves $<,\le,=,\ge,>$. The left-hand side is computed exactly by \cref{lem:overflow}.
\end{proof}

\begin{theorem}[Kernel correctness]\label{thm:kernel}
Let $M$ be well-formed and initialisable, with $\hat M=\code{Model::create}(M)$.
Then $(\abs^{-1},\mathrm{id}_T,\mathrm{id}_P)$ is an isomorphism
$\sem{M}\iso\Ksem(M)$, where an IR delay $d$ corresponds to the kernel delay $dR$.
In detail, for all $s\in\Valid(M)$ and $\kappa=\abs^{-1}(s)$:
\begin{enumerate}[label=(\roman*)]
\item \textbf{Initial state.} \code{initial\_configuration} returns $\ok(\abs^{-1}(s_0))$.
\item \textbf{Delay soundness and completeness.} $s\trans{d}s'$ in $\sem{M}$ iff
  $\code{delay}(\hat M,\kappa,dR)=\ok(\abs^{-1}(s'))$.
\item \textbf{Discrete soundness and completeness.} $s\trans{t_j}s'$ in $\sem{M}$ iff
  $\code{fire}(\hat M,\kappa,j)=\ok(\abs^{-1}(s'))$.
\item \textbf{Propositions.} $p\in\lab_M(s)$ iff $p\in\code{propositions}(\hat M,\kappa)$.
\end{enumerate}
\emph{Soundness} (the kernel cannot invent behaviour) is the right-to-left direction
of (ii) and (iii). \emph{Completeness} (the kernel drops no behaviour) is the
left-to-right direction.
\end{theorem}
\begin{proof}
(i) The kernel returns $(\iota,0^n,0)=\abs^{-1}(\iota,\mathbf{0},0)$ iff that configuration
satisfies $\hat J_\iota$, which by \cref{lem:atom} holds iff $\mathbf{0}\sat J_\iota$, the
premise of \textsc{(Init$_M$)}.

(ii) Let $s=(l,w,\theta)$ and $\kappa=(l,c,\tau)$, so $c=Rw$ and $\tau=R\theta$. The kernel
returns $\ok(\kappa')$ exactly when all of the following hold: $dR\ge0$; $\kappa$
satisfies $\hat J_l$; $\tau+dR\le\Tmax$ and $c_i+dR\le\Tmax$ for all $i$; and
$(l,c+dR,\tau+dR)$ satisfies $\hat J_l$. In that case $\kappa'=(l,c+dR,\tau+dR)$. Since
$c_i\le\tau$, the clock bounds follow from the time bound, which in turn is
$\theta+d\le\Hz$. By \cref{lem:atom} the two invariant checks are $w\sat J_l$ and
$w+d\sat J_l$. These are exactly the premises of \textsc{(Delay$_M$)}, and
$\abs(\kappa')=(l,w+d,\theta+d)$ is its target. If the premises fail, the kernel returns
an error, and $s$ has no $d$-successor in $\sem{M}$.

(iii) Let $t_j=(\cdot,\mathit{src},\mathit{tgt},\alpha,g,Y)$. The kernel returns $\ok$ iff
$l=\mathit{src}$, $\kappa$ satisfies $\hat g$, and $(\mathit{tgt},c[Y{:=}0],\tau)$ satisfies
$\hat J_{\mathit{tgt}}$; the result is that configuration. By \cref{lem:atom}, and
because $\abs$ commutes with resets ($\abs(l,c[Y{:=}0],\tau)$ has valuation $w[Y{:=}0]$),
these are the premises of \textsc{(Disc$_M$)} and the targets coincide. The index
range check never fails for $j<m$.

(iv) Both labellings are $\{p_l\}$.

The maps $\abs^{-1}$, $\mathrm{id}_T$ and $\mathrm{id}_P$ are bijections
(\cref{lem:alpha}), so (i)--(iv) are the conditions of \cref{def:iso}.
\end{proof}

\subsection{Derived kernel services}

\begin{lemma}[Enabling windows are exact]\label{lem:window}
For valid $\kappa$ and transition $t_j$, \code{enabling\_window} returns an interval
$[a,b]$ (with $b=\infty$ meaning unbounded up to the horizon) or nothing, and
\[
  \delta\in[a,b]\cap\Nat \iff \code{fire}(\hat M,\code{delay}(\hat M,\kappa,\delta),j)\text{ is }\ok .
\]
\end{lemma}
\begin{proof}
Every check made by \code{delay} and \code{fire} after a delay $\delta$ is an atom
evaluated at a valuation $u+\delta s$, where $u$ is the current (or post-reset)
valuation and $s_x\in\{0,1\}$ (0 for the reference clock and for reset clocks).
Each atom becomes $d_0+k\delta\bowtie\hat c$ with $k\in\{-1,0,1\}$. For $k=0$ this is a
constant truth value. For $k=\pm1$ it is a half-line or a point of integers, which
\code{restrict\_delay} intersects exactly, translating strict bounds by $\pm1$
because $\delta$ ranges over the integers. The horizon contributes
$\delta\le\Tmax-\tau$ and $\delta\le\Tmax-c_i$. The intersection of all these sets is
exactly the set of $\delta$ for which every check passes.
\end{proof}

\begin{definition}[Monitoring semantics]\label{def:monitor}
For a non-empty finite set $S$ of kernel configurations sharing the time $\tau$,
an observation $(\tau',\sigma)$, where $\sigma$ is an action label or a transition
index, yields
\[
  \mathrm{obs}(S,\tau',\sigma)=\{\kappa'\mid \exists\kappa\in S,\,j\text{ matching }\sigma:\
    \code{fire}(\code{delay}(\kappa,\tau'-\tau),j)=\ok(\kappa')\},
\]
and is rejected if this set is empty or $\tau'<\tau$. This is \code{kernel::observe}.
\end{definition}

\begin{proposition}[Monitor exactness]\label{prop:monitor}
After observations $o_1\cdots o_r$ starting from $\{\kappa_0\}$, the monitor's
set is exactly
$\{\abs^{-1}(s)\mid s_0\text{ reaches }s\text{ in }\sem{M}\text{ by a run whose discrete
steps match }o_1,\dots,o_r\text{ at the observed times}\}$, where intermediate
internal steps are themselves observations, as in \cref{rem:tau}. In particular,
for event-deterministic models, the set is always a singleton and the monitor
coincides with $\Ksem$.
\end{proposition}
\begin{proof}
By induction on $r$, using \cref{thm:kernel} for each delay/discrete pair and
the definition of $\mathrm{obs}$ as the image of the step relation. Pure delay
observations (\code{advance\_to}) are handled the same way.
\end{proof}

\begin{remark}[Internal steps]\label{rem:tau}
$\tauact$-transitions of $\VD$ are internal decisions of the twin itself. The
runtime fires them explicitly, by transition identifier and at a logical time,
so the kernel never has to guess when they happened. The monitor therefore
never computes $\tauact$-closures over unknown delays, which would require
symbolic zones in the trusted kernel.
\end{remark}

\begin{proposition}[Prediction does not interfere]\label{prop:nointerf}
No prediction or query function (\code{enabled\_transitions}, \code{successors},
\code{enabling\_window}, \code{explore}, \code{simulate}, \code{evaluate\_propositions})
changes the authoritative state. Snapshot restoration accepts exactly the
admissible configurations $\W(\hat M)$.
\end{proposition}
\begin{proof}
The free functions take their state arguments by \code{const} reference and
return fresh values. The corresponding \code{KernelInstance} methods are
\code{const} member functions, so the C++ type system forbids mutating the
held state. \code{simulate} copies the start state explicitly. The mutators
(\code{advance\_*}, \code{apply\_*}, \code{restore\_state}) compute the candidate state
first and assign it only on success. \code{restore\_state} applies
\code{check\_configuration} to every member, and that function implements
membership in $\W(\hat M)$ (\cref{def:kconf}).
\end{proof}

% ===========================================================================
\section{The production runtime: $\Ksem\wbis\Kprod$}\label{sec:prod}
% ===========================================================================

The production runtime (\code{twin-runtime}, library \code{twin::runtime}) contains
the kernel together with package loading, input adapters, an event sequencer,
the execution ledger, REST/SSE APIs, metrics, snapshots, the mission
controller and the co-simulation driver. We model it at the level of
granularity at which it interacts with semantics.

\subsection{Semantic and internal actions}

\begin{definition}[Production LTTS $\Kprod$]\label{def:kprod}
The states of $\Kprod(M)$ are pairs $(\kappa,b)$ with $\kappa\in\W(\hat M)$, the
authoritative kernel configuration held by the session, and
$b\in\mathcal{B}$, the \emph{bookkeeping state}: ledger file contents, event
buffers, API caches, metrics, snapshot files, adapter queues, planner caches,
and so on. The labels are the kernel labels (transitions $j$ and delays $\delta$)
plus a set $\Sigma_{\mathrm{int}}$ of internal runtime actions, all of which are
internal ($\tauact$):
\begin{center}
\begin{tabular}{@{}ll@{}}
\toprule
Internal action ($\tauact$) & Effect \\ \midrule
ledger append / fsync / anchor & $b$ only \\
API request parsing and response serialisation & $b$ only \\
SSE / WebSocket publication & $b$ only \\
metrics and structured logging & $b$ only \\
snapshot persistence & $b$ only \\
adapter polling (telemetry, map updates), queueing & $b$ only \\
planning, plan validation, prediction queries & $b$ only (\cref{prop:nointerf}) \\
\bottomrule
\end{tabular}
\end{center}
The transitions are
\begin{align*}
(\kappa,b)&\trans{\lambda}(\kappa',b') && \text{if } \kappa\trans{\lambda}\kappa' \text{ in }\Ksem
   \text{ and } b'=\mathit{rec}_\lambda(b,\kappa,\kappa'), \quad \lambda\in\{0..m-1\}\cup\Nat,\\
(\kappa,b)&\trans{\tauact}(\kappa,b') && \text{for every internal action with effect } b\mapsto b',
\end{align*}
with labelling $\lab(\kappa,b)=\lab_K(\kappa)$. Here $\mathit{rec}$ is the
bookkeeping that accompanies a committed step.
\end{definition}

The definition encodes three architectural properties. The theorem
depends on them, and the implementation is built to ensure them.

\begin{assumption}[Semantic isolation]\label{ass:isolation}
No internal action modifies $\kappa$. In the code, the authoritative state is
the private \code{KernelInstance} inside \code{TwinSession}. The session exposes
only const views to adapters, the API, the ledger writer, the planner and the
UI. The architecture tests check that runtime sources outside the session do
not construct or assign kernel configurations.
\end{assumption}

\begin{assumption}[Single semantic authority]\label{ass:authority}
Every change of $\kappa$ is the result of a kernel transition function
(\code{delay}/\code{fire} via \code{KernelInstance}). There is no second
implementation of the transition relation in the shell, adapters, planner or UI.
The architecture tests check that runtime, planner and UI sources contain no
location names of the model, so no shadow state machine can branch on them.
\end{assumption}

\begin{assumption}[Availability: no semantic refusal by the shell]\label{ass:available}
Every well-formed request for a kernel step is forwarded to the kernel. Between
two requests, the runtime performs only finitely many internal actions. The
shell refuses a well-formed request only when the kernel refuses it, or when
infrastructure fails (for example, the ledger cannot be written), in which case
the runtime \emph{fails stop}: it accepts no further semantic steps.
\end{assumption}

\subsection{The weak bisimulation}

\begin{theorem}[Production correctness]\label{thm:production}
Under \cref{ass:isolation,ass:authority,ass:available}, and in the absence of
infrastructure failures, the relation
\[
  \mathcal{R}=\{\,(\kappa,(\kappa,b))\mid \kappa\in\W(\hat M),\ b\in\mathcal{B}\,\}
\]
is a weak timed bisimulation between $\Ksem(M)$ and $\Kprod(M)$ that relates
the initial states. Hence $\Ksem(M)\wbis\Kprod(M)$.
\end{theorem}
\begin{proof}
The initial states are related because the runtime initialises the session
with \code{KernelInstance::initialize}, that is, with \code{initial\_configuration}.
Let $(\kappa,(\kappa,b))\in\mathcal{R}$.

(W1) $\lab(\kappa,b)=\lab_K(\kappa)$ by definition.

(W2) Suppose $\kappa\trans{\lambda}\kappa'$ in $\Ksem$, with $\lambda$ a transition or a
delay. By \cref{ass:available}, a request for $\lambda$ reaches the kernel after
finitely many internal actions $(\kappa,b)\wtrans{\epsilon}(\kappa,b_1)$; these leave
$\kappa$ unchanged by \cref{ass:isolation}. The kernel computes the same function
as in $\Ksem$, so $(\kappa,b_1)\trans{\lambda}(\kappa',b_2)$ with
$b_2=\mathit{rec}_\lambda(b_1,\kappa,\kappa')$, and $(\kappa',(\kappa',b_2))\in\mathcal{R}$.
If $\lambda$ is internal in $\Ksem$ (a $\tauact$-transition of the model), the same
step is matched by a $\tauact$-step of $\Kprod$, which is admissible under
$\wtrans{\epsilon}$.

(W3) Suppose $(\kappa,b)\trans{\mu}(\kappa'',b'')$ in $\Kprod$. If $\mu$ is an internal
runtime action, then $\kappa''=\kappa$ by \cref{ass:isolation}, and $\kappa$ matches it
by the empty move, with $(\kappa,(\kappa,b''))\in\mathcal{R}$. Otherwise, by
\cref{def:kprod} and \cref{ass:authority}, $\mu=\lambda$ is a kernel step
$\kappa\trans{\lambda}\kappa''$, which $\Ksem$ matches exactly. Delays of $\Kprod$ are
kernel delays, because logical time is a component of $\kappa$ and advances only
through \code{delay}, so delay labels are matched exactly as well.
\end{proof}

\begin{corollary}[Safety under infrastructure failures]\label{cor:failstop}
Without \cref{ass:available}, but with \cref{ass:isolation,ass:authority}, the
relation $\mathcal{R}^{-1}$ is still a weak timed \emph{simulation} of $\Kprod$ by
$\Ksem$. Every observable timed trace of the runtime, including traces cut short
by a fail-stop, is a trace of $\Ksem$. Infrastructure failures can stop the twin;
they cannot make it produce behaviour that the model does not admit.
\end{corollary}
\begin{proof}
Clause (W3) of the proof above uses only \cref{ass:isolation,ass:authority}.
\end{proof}

\begin{remark}[The ledger is an internal action]\label{rem:ledger}
A committed step is recorded in the ledger before the session publishes the new
state (write-ahead), and the hash chain is computed from the committed values.
The ledger therefore records the run, but it never influences which
transitions exist. \Cref{thm:production} is independent of the ledger content,
and replay (re-executing the recorded inputs with the kernel) reproduces the
recorded states by determinism of the kernel functions.
\end{remark}

\begin{remark}[Monitor sets and execution policies]\label{rem:policy}
For models that are not event-deterministic, the session holds the monitor set of
\cref{def:monitor} instead of a single $\kappa$. \Cref{thm:production} then holds
verbatim with $\kappa$ replaced by that set and $\Ksem$ replaced by its monitor
(subset) construction, which is trace-equivalent to $\Ksem$ by \cref{prop:monitor}.
Where an execution must commit to one successor (for example, to command the
physical twin), the choice is made by an explicit \emph{policy} outside the
kernel. A policy-constrained execution is a refinement: a weak timed
\emph{simulation} by $\Ksem$, not a bisimulation. The demonstration model is
event-deterministic (checked by the compiler), so neither situation arises there.
\end{remark}

% ===========================================================================
\section{Composition and the alignment chain}\label{sec:composition}
% ===========================================================================

\begin{theorem}[$\VD$ is realised by the runtime]\label{thm:composition}
Let $\VD$ be a UPPAAL document whose compilation succeeds, with $M=\IR(\VD)$.
Under \cref{ass:utap,ass:isolation,ass:authority,ass:available} and the
logical-time semantics of \cref{def:grid},
\[
  \sem{\VD}\ \iso\ \sem{M}\ \iso\ \Ksem(M)\ \wbis\ \Kprod(M),
  \qquad\text{hence}\qquad \sem{\VD}\ \wbis\ \Kprod(M).
\]
The labels are identified by the action-preserving maps $\nu_\Edges$ and
$\mathrm{id}$, the propositions by $\at{\ell}\mapsto p_{\nu_\Loc(\ell)}$, and delays by
$d\mapsto dR$ ticks.
\end{theorem}
\begin{proof}
\Cref{thm:compiler,thm:kernel,thm:production}, followed by \cref{lem:facts}(a)
to weaken isomorphisms to weak timed bisimulations and \cref{lem:facts}(b) to
compose. The label and proposition maps compose because each step preserves
actions and propositions.
\end{proof}

\begin{corollary}[Grid-timed traces]\label{cor:traces}
Every observable timed trace of the deployed twin $D=\Kprod(\IR(\VD))$ is a timed
trace of the dense-time semantics $\sem{\VD}_{\Real}$. Every timed trace of
$\sem{\VD}_{\Real}$ whose timestamps lie on the grid $\TR^{\Hz}$ is a timed trace of $D$
(under \cref{ass:available}).
\end{corollary}
\begin{proof}
Combine \cref{thm:composition} (weak bisimilarity implies equality of observable
timed traces) with \cref{prop:dense}.
\end{proof}

\subsection{Connection to semantic alignment}

The alignment paper's chain $P\simul\VP\sbis_{\cons}\VD\sbis D$ has three links, and
this project discharges only the last one:

\begin{center}
\begin{tabularx}{\linewidth}{@{}lXl@{}}
\toprule
Link & Status in this project & Evidence \\ \midrule
$P\simul\VP$ & assumption about the physical system and the PT engineering process & none (out of scope) \\
$\VP\sbis_{\cons}\VD$ & assumption, decided by the unmodified aligner (Z3 + zone graphs) &
  \code{evidence/alignment.json} in the package, hash-bound to $\VP$, $\VD$, $K$, $I_P$, $I_D$ \\
$\VD\sbis D$ & \emph{theorem} (\cref{thm:composition}), in its weak-bisimulation form over $\TR^{\Hz}$ &
  this report, plus compilation manifest and translation validation \\
\bottomrule
\end{tabularx}
\end{center}

Because $\VD\wbis D$ preserves observable propositions and timed traces, the
guarantees the alignment paper derives from $\VP\sbis_{\cons}\VD$ apply to $D$
within the logical-time semantics. These are agreement of the twins' domain
interpretations at matched states (paper, Theorem~1) and transfer of TCTL
properties over the observable propositions (paper, Theorem~2). Two
qualifications must be stated:
\begin{enumerate}
\item The guarantee is only as strong as the alignment verdict. The integration
  report lists the aspects of Definition~5 that the aligner's implementation
  does not check (Condition~I is not evaluated; uninterpreted synchronised
  labels are skipped; the zone graph starts from the unconstrained initial
  zone). These qualify the \emph{assumption} $\VP\sbis_{\cons}\VD$, not \cref{thm:composition}.
\item Properties of $\sem{\VD}_{\Real}$ that depend on off-grid timing need not hold
  for $D$, and vice versa. By \cref{cor:traces}, universally quantified trace
  properties (safety) of $\sem{\VD}_{\Real}$ hold for $D$.
\end{enumerate}

\paragraph{The defensible product statement.}
\emph{A verified DT view is compiled into a canonical executable representation
and executed by a small trusted semantic kernel whose logical behaviour realises
the verified view by construction (\cref{thm:composition}). Runtime decisions and
observations are recorded in a cryptographically tamper-evident execution
ledger, while planners and industrial integrations remain outside the trusted
semantic core (\cref{ass:isolation,ass:authority}).}

% ===========================================================================
\section{Logical time in the runtime}\label{sec:time}
% ===========================================================================

\paragraph{Logical, not physical.}
All time in the theorems is \emph{model time}. A timestamp $t$ on an observation
means ``in the model, this event happens when the global logical clock $\theta$
equals $t$''. The runtime does not read the operating system clock in order to
decide semantic questions. The wall-clock time appears only as optional,
informational metadata in ledger records and is never consulted by the kernel.
The theorems therefore say nothing about how logical time relates to physical
time; establishing that relation (sensor timestamping, clock synchronisation,
bounded processing latency) is outside the scope of this report
(\cref{sec:limitations}).

\paragraph{Exact grid.}
Timestamps are parsed exactly from decimal strings (\code{twin::parse\_time}):
\code{"31.5"} at $R=1000$ is $31500$ ticks. A timestamp with more fractional
digits than the grid allows yields the error \code{TimeNotRepresentable}; it is
never rounded. The APIs reject JSON floating-point numbers for times.

\paragraph{Observation processing.}
An observation $(t,\sigma)$ is processed as a timed step: delay by $t-\theta$, then
the discrete step selected by $\sigma$ (\cref{def:monitor}). Hence:
\begin{itemize}
\item if $t<\theta$, the observation is rejected (\code{TimeRegression});
\item if the delay violates the current invariant, the observation is rejected
  (\code{IncompatibleObservation}, reporting the deadline \code{max\_delay}). For
  example, if the twin is in a location with invariant $x\le 5$ and $x=1.5$ at
  $\theta=28$, an event at $t=31.5$ requires $x=5$ and is admissible, whereas an
  event at $t=31.6$ is rejected;
\item a rejected observation leaves the semantic state unchanged and is recorded
  in the ledger as a rejection.
\end{itemize}

\paragraph{Simultaneous events.}
Events carrying the same timestamp are processed one after the other with zero
delay between them, in the order fixed by the runtime's single sequencer: the
order in which the input adapters enqueue them, with a deterministic
tie-breaking order between adapters. The ledger records that order, and replay
follows it. This is a total order on inputs; it is part of the input, not of
the model.

\paragraph{Pure passage of time.}
Between events the runtime may let time pass without an event
(\code{advance\_to}). This is the \textsc{(Delay)} rule. If the passage of time
alone would violate the current invariant (a missed deadline), the delay is
refused, the runtime reports a monitoring alarm, and the state is not changed.

% ===========================================================================
\section{Assumptions, trusted base and limitations}\label{sec:limitations}
% ===========================================================================

\subsection{What is established}
Under the stated assumptions, the relationship between the formal DT view
$\VD$, its compiled IR, and the logical semantic execution performed by the
kernel and the production runtime is an isomorphism, respectively a weak
timed bisimulation, over the logical-time grid $\TR^{\Hz}$ (\cref{thm:composition}).

\subsection{Assumptions}
\begin{enumerate}[label=(A\arabic*)]
\item \textbf{Source reading}: UTAP parses UPPAAL XML as specified (\cref{ass:utap}).
  The aligner rests on the same assumption.
\item \textbf{Implementation = definitions}: the C++ functions listed in
  \cref{app:code} compute the mathematical functions of
  \cref{sec:ir,sec:compiler,sec:kernel}. The kernel was kept small (about 600
  lines including documentation) to make this assumption auditable, and the
  compiler's output is translation-validated (\cref{prop:tv}).
\item \textbf{C++ toolchain and platform}: the compiler, the standard library and
  the hardware execute the code as the C++ standard prescribes. In particular,
  \code{int64} arithmetic is two's complement and \code{\_\_builtin\_add\_overflow}
  reports overflow.
\item \textbf{Runtime architecture}: \cref{ass:isolation,ass:authority} (enforced by
  code structure and architecture tests) and \cref{ass:available} (availability;
  without it only \cref{cor:failstop} holds).
\item \textbf{Alignment}: $\VP\sbis_{\cons}\VD$ is assumed, and the evidence is the
  verdict of the unmodified aligner, qualified by the integration findings.
\end{enumerate}

\subsection{Trusted computing base for $\VD\wbis D$}
\begin{center}
\begin{tabularx}{\linewidth}{@{}lX@{}}
\toprule
Must be trusted & Why \\ \midrule
\code{twin::kernel} & computes every semantic successor (\cref{thm:kernel}) \\
\code{twin::ir\_model} (types, \code{validate}) & defines what the kernel executes \\
\code{twin::core} (\code{Result}, logical time) & exact time arithmetic and parsing \\
IR loader (\code{ir::from\_canonical\_text}) + SHA-256 & loads exactly the hashed IR (\cref{lem:canon}) \\
\code{twin::compiler} \emph{or} its translation validation & establishes $\VD\iso\IR(\VD)$ \\
UTAP (shared with the aligner) & reads the source \\
\code{TwinSession} commit protocol & the single mutation path (\cref{ass:authority}) \\
\bottomrule
\end{tabularx}
\end{center}

\noindent\emph{Need not be trusted for $\VD\wbis D$.} These components may fail
without invalidating semantic correctness; they can at worst cause refusals,
stops or wrong \emph{proposals}. They are the planner, the plan validator, the
mission controller, the input adapters (they produce candidate observations,
and wrong observations are rejected or faithfully recorded), the map service,
the PT simulator, the visualisation, the HTTP/SSE server, metrics and logging,
and the ledger (it records the run but cannot change it; it matters for
\emph{audit}, not for semantics).

\noindent\emph{What would invalidate the claim.} Any of the following would: a
code path that changes the session's configuration without a kernel call; a
second implementation of transition logic in the shell, planner or UI; a kernel
or IR-validation defect; loading an IR other than the hashed one; changing
$\VD$ after the alignment verdict without re-running the aligner (the package
hashes detect this); or using a UTAP version that parses differently from the
one the alignment was performed with.

\subsection{What is \emph{not} established}
\begin{itemize}
\item hard real-time execution, worst-case execution time, or response-time bounds;
\item equivalence between logical timestamps and wall-clock time;
\item operating-system scheduler, network latency or message-ordering guarantees;
\item synchronisation of distributed clocks;
\item equivalence between the physical plant and any model ($P\simul\VP$ is assumed);
\item correctness, optimality or completeness of the path planner;
\item correctness of the building sensor model or of the environment service;
\item the soundness of the alignment verdict beyond what the aligner checks
  (integration findings 10--13);
\item mechanised (machine-checked) correctness of this proof. Mechanisation in
  Lean, Coq or Isabelle/HOL is future work.
\end{itemize}

\subsection{Future work}
Mechanising \cref{thm:compiler,thm:kernel} is the natural next step: their
proofs are by direct case analysis over a small rule set, which makes them
well suited to a proof assistant, possibly with code extraction for the
kernel. A second direction is a hard-real-time realisation that relates
logical time to physical time under explicit WCET and scheduling assumptions.

\appendix
% ===========================================================================
\section{Correspondence between definitions and code}\label{app:code}
% ===========================================================================

\begin{center}
\small
\begin{tabularx}{\linewidth}{@{}lX@{}}
\toprule
Formal object & Implementation \\ \midrule
\cref{def:grid} (grid $\TR$, horizon) & \code{twin::TimeBase}, \code{kMaxTicks}, \code{parse\_time}, \code{format\_time} (\code{src/core/logical\_time.cpp}) \\
\cref{def:source} (source automaton) & \code{compiler::detail::SourceModel} (\code{src/compiler/source\_model.hpp}) \\
\cref{rem:side} (side conditions) & \code{StrictReader} (\code{src/compiler/utap\_reader.cpp}), diagnostics \code{TWC0xx} \\
\cref{def:ir,def:wf-ir} (IR, well-formedness) & \code{twin::ir::Model}, \code{ir::validate} (\code{src/ir/validate.cpp}) \\
\cref{def:canon} (canonical text) & \code{ir::to\_canonical\_text}, \code{json::canonical\_dump} \\
\cref{lem:canon} (faithful loading) & \code{ir::from\_canonical\_text} (round-trip check) \\
\cref{def:compiler} (compiler $C$) & \code{compiler::compile\_file} (\code{src/compiler/compiler.cpp}) \\
\cref{prop:tv} (translation validation) & \code{detail::crosscheck\_with\_aligner} (\code{src/compiler/crosscheck.cpp}) \\
\cref{def:kmodel} (kernel model) & \code{kernel::Model::create} (\code{src/kernel/model.cpp}) \\
\cref{def:kconf} (configurations, $\W$) & \code{kernel::Configuration}, \code{check\_configuration} \\
\cref{def:katom} (atom satisfaction) & \code{atom\_holds}, \code{satisfies} (\code{src/kernel/semantics.cpp}) \\
\cref{def:kfun} (transition functions) & \code{initial\_configuration}, \code{delay}, \code{fire}, \code{propositions} \\
\cref{lem:window} (enabling windows) & \code{enabling\_window}, \code{restrict\_by\_atom} \\
\cref{def:monitor} (monitoring) & \code{kernel::observe}, \code{advance\_to} (\code{src/kernel/state\_set.cpp}) \\
\cref{prop:nointerf} (prediction) & \code{explore}, \code{simulate}, \code{KernelInstance} const members \\
\cref{def:kprod} ($\Kprod$) & \code{runtime::TwinSession} and the shell modules (\code{src/runtime}) \\
\cref{rem:ledger} (ledger) & \code{twin::ledger} (\code{src/ledger}) \\
\bottomrule
\end{tabularx}
\end{center}

\noindent The test suites cross-check several of these correspondences
empirically: compiler determinism and translation validation on the aligner's
corpus, kernel-versus-aligner zone-graph agreement on generated automata, and
exhaustive window/step agreement around boundaries. These tests are validation
and do not replace the proofs.


\begin{thebibliography}{9}
\bibitem{alur-dill} R.~Alur and D.~L.~Dill. A theory of timed automata.
  \emph{Theoretical Computer Science}, 126(2):183--235, 1994.
\bibitem{uppaal} G.~Behrmann, A.~David, K.~G.~Larsen. A tutorial on \textsc{Uppaal}.
  In \emph{SFM-RT}, LNCS 3185, 2004.
\bibitem{hmp} T.~A.~Henzinger, Z.~Manna, A.~Pnueli. What good are digital clocks?
  In \emph{ICALP}, LNCS 623, 1992.
\bibitem{milner} R.~Milner. \emph{Communication and Concurrency}. Prentice Hall, 1989.
\bibitem{cerans} K.~\v{C}er\=ans. Decidability of bisimulation equivalences for
  parallel timer processes. In \emph{CAV}, LNCS 663, 1992.
\bibitem{tv} A.~Pnueli, M.~Siegel, E.~Singerman. Translation validation.
  In \emph{TACAS}, LNCS 1384, 1998.
\bibitem{icse} A Formal Software Engineering Methodology for Semantically Aligned
  Digital Twins. Paper accompanying \emph{SemPTDTAlignmentICSE} (cited as ``the
  alignment paper'').
\bibitem{fips} NIST. FIPS PUB 180-4: Secure Hash Standard, 2015.
\bibitem{jcs} A.~Rundgren, B.~Jordan, S.~Erdtman. RFC 8785: JSON Canonicalization Scheme, 2020.
\end{thebibliography}

\end{document}
